% ============================================================
%  आवरण पृष्ठ (Front cover) — बाल-सुलभ, रंगीन, चित्रमय डिज़ाइन
% ============================================================
% ============================================================
%  cover-graphics.tex — साझा रंग, फ़ॉन्ट एवं TikZ चित्र
%  (आवरण front + back दोनों इसे \input करते हैं — guard ज़रूरी)
% ============================================================
\ifcsname cvGraphicsLoaded\endcsname\else
\expandafter\gdef\csname cvGraphicsLoaded\endcsname{}

% ---------------- बाल-सुलभ रंग-पट्टिका (playful palette - kept for compatibility) ----------------
\definecolor{cvSkyTop}{HTML}{63C7F2}
\definecolor{cvSkyBot}{HTML}{EAF8FF}
\definecolor{cvHill}{HTML}{7FC85A}
\definecolor{cvHillD}{HTML}{67B048}
\definecolor{cvBlue}{HTML}{1F4E79}
\definecolor{cvOrange}{HTML}{F5821F}
\definecolor{cvYellow}{HTML}{FFC83D}
\definecolor{cvRed}{HTML}{E5524A}
\definecolor{cvPurple}{HTML}{8E5BD9}
\definecolor{cvCyan}{HTML}{2FB4D6}
\definecolor{cvPink}{HTML}{FF7BAC}
\definecolor{cvLime}{HTML}{9CD356}
\definecolor{cvGrey}{HTML}{5F6368}

% ---------------- मॉडर्न विज्ञान/तकनीक रंग-पट्टिका (Modern Tech Palette) ----------------
\definecolor{cvBgDark}{HTML}{080D1A}
\definecolor{cvBgMid}{HTML}{0E1830}
\definecolor{cvBgLight}{HTML}{1A2B4C}
\definecolor{cvNeonCyan}{HTML}{00F2FE}
\definecolor{cvNeonBlue}{HTML}{4FACFE}
\definecolor{cvNeonPurple}{HTML}{E255FF}
\definecolor{cvNeonPink}{HTML}{FF0844}
\definecolor{cvNeonGreen}{HTML}{38EF7D}
\definecolor{cvNeonYellow}{HTML}{F9D423}
\definecolor{cvNeonOrange}{HTML}{FF5E62}
\definecolor{cvGridColor}{HTML}{14244A}

% ---------------- फ़ॉन्ट (Baloo 2 — rounded, Devanagari + Latin) ----------------
% ucharclasses को स्थानीय रूप से बंद कर (\XeTeXinterchartokenstate=0) एक ही
% फ़ॉन्ट से दोनों लिपियाँ — एकरूप, मज़ेदार शैली।
% Devanagari shaping needs Script=Devanagari; Latin uses default. अलग families।
\newfontfamily\cvTitleDev{Baloo2-ExtraBold.ttf}[Path=fonts/, Script=Devanagari, Language=Hindi]
\newfontfamily\cvHeadDev{Baloo2-SemiBold.ttf}[Path=fonts/, Script=Devanagari, Language=Hindi]
\newfontfamily\cvBodyDev{Baloo2-Medium.ttf}[Path=fonts/, Script=Devanagari, Language=Hindi]
\newfontfamily\cvTitleLat{Baloo2-ExtraBold.ttf}[Path=fonts/]
\newfontfamily\cvHeadLat{Baloo2-SemiBold.ttf}[Path=fonts/]
\newfontfamily\cvBodyLat{Baloo2-Medium.ttf}[Path=fonts/]
% \cvXX : ucharclasses बंद कर Baloo चुनें — H=हिन्दी, L=Latin
\newcommand{\cvTH}{\XeTeXinterchartokenstate=0 \cvTitleDev}
\newcommand{\cvTL}{\XeTeXinterchartokenstate=0 \cvTitleLat}
\newcommand{\cvHH}{\XeTeXinterchartokenstate=0 \cvHeadDev}
\newcommand{\cvHL}{\XeTeXinterchartokenstate=0 \cvHeadLat}
\newcommand{\cvBH}{\XeTeXinterchartokenstate=0 \cvBodyDev}
\newcommand{\cvBL}{\XeTeXinterchartokenstate=0 \cvBodyLat}

% ---------------- चित्र (pics) ----------------
\tikzset{
  planet/.pic={
    \draw[cvNeonPurple, line width=2pt] (0,0) ellipse (1.15 and 0.42);
    \fill[cvNeonOrange] (0,0) circle (0.7);
    \fill[cvNeonOrange!70] (0.22,0.18) circle (0.18);
    \fill[cvNeonOrange!70!black] (-0.25,-0.12) circle (0.12);
    \draw[cvNeonPurple, line width=2pt] (1.15,0) arc (0:180:1.15 and 0.42);
  },
  rocket/.pic={
    \fill[cvNeonOrange] (0,-0.95) -- (-0.2,-0.5) -- (0,-0.64) -- (0.2,-0.5) -- cycle;
    \fill[cvNeonYellow] (0,-0.8) -- (-0.1,-0.5) -- (0,-0.6) -- (0.1,-0.5) -- cycle;
    \fill[cvNeonPink] (-0.3,-0.5) -- (-0.56,-0.78) -- (-0.3,-0.18) -- cycle;
    \fill[cvNeonPink] (0.3,-0.5) -- (0.56,-0.78) -- (0.3,-0.18) -- cycle;
    \fill[white] (-0.3,-0.5) .. controls (-0.32,0.45) .. (0,1.0)
                 .. controls (0.32,0.45) .. (0.3,-0.5) -- cycle;
    \draw[cvNeonBlue, line width=1.5pt] (-0.3,-0.5) .. controls (-0.32,0.45) .. (0,1.0)
                 .. controls (0.32,0.45) .. (0.3,-0.5) -- cycle;
    \fill[cvNeonPink] (0,1.0) .. controls (0.18,0.55) .. (0.15,0.35) -- (-0.15,0.35)
                 .. controls (-0.18,0.55) .. cycle;
    \fill[cvNeonCyan] (0,0.05) circle (0.15);
    \draw[cvNeonBlue, line width=1pt] (0,0.05) circle (0.15);
  },
  atom/.pic={
    \draw[cvNeonCyan, line width=1.4pt] (0,0) ellipse (0.95 and 0.36);
    \draw[cvNeonCyan, line width=1.4pt, rotate=60] (0,0) ellipse (0.95 and 0.36);
    \draw[cvNeonCyan, line width=1.4pt, rotate=-60] (0,0) ellipse (0.95 and 0.36);
    \fill[cvNeonPink] (0,0) circle (0.17);
    \fill[cvNeonBlue] (0.95,0) circle (0.08);
    \fill[cvNeonBlue] (-0.48,0.82) circle (0.08);
    \fill[cvNeonBlue] (-0.48,-0.82) circle (0.08);
  },
  dna/.pic={
    \draw[cvNeonPink, line width=1.8pt] plot[domain=0:360, samples=46] ({0.32*sin(\x)},{\x/130 - 1.38});
    \draw[cvNeonPurple, line width=1.8pt] plot[domain=0:360, samples=46] ({-0.32*sin(\x)},{\x/130 - 1.38});
    \foreach \t in {25,65,...,335}
      \draw[cvNeonCyan, line width=1pt] ({0.32*sin(\t)},{\t/130 - 1.38}) -- ({-0.32*sin(\t)},{\t/130 - 1.38});
  },
  flask/.pic={
    \fill[cvNeonCyan!20] (-0.1,0.6) -- (-0.1,0.16) -- (-0.45,-0.5) -- (0.45,-0.5) -- (0.1,0.16) -- (0.1,0.6) -- cycle;
    \fill[cvNeonGreen] (-0.3,-0.15) -- (-0.45,-0.5) -- (0.45,-0.5) -- (0.3,-0.15) -- cycle;
    \draw[cvNeonBlue, line width=1.3pt] (-0.1,0.6) -- (-0.1,0.16) -- (-0.45,-0.5) -- (0.45,-0.5) -- (0.1,0.16) -- (0.1,0.6);
    \draw[cvNeonBlue, line width=1.3pt] (-0.17,0.6) -- (0.17,0.6);
    \fill[white] (0.07,-0.3) circle (0.045);
    \fill[white] (-0.12,-0.38) circle (0.03);
  },
  bulb/.pic={
    \fill[cvNeonYellow] (0,0) circle (0.33);
    \draw[cvNeonOrange, line width=1pt] (0,0) circle (0.33);
    \fill[cvGrey!55] (-0.14,-0.3) rectangle (0.14,-0.46);
    \draw[cvNeonOrange, line width=1pt] (-0.08,0.12) -- (0,-0.04) -- (0.08,0.12);
    \foreach \a in {45,90,135} \draw[cvNeonYellow, line width=1.4pt] (\a:0.42) -- (\a:0.58);
  },
  robot/.pic={
    \draw[cvGrey, line width=2.4pt] (0,1.0) -- (0,1.35);
    \fill[cvRed] (0,1.45) circle (0.12);
    \fill[white] (-1,-0.85) rectangle (1,1.0);
    \draw[cvBlue, line width=2.4pt, rounded corners=10pt] (-1,-0.85) rectangle (1,1.0);
    \fill[cvRed, rounded corners=2pt] (-1.2,0.0) rectangle (-1.0,0.45);
    \fill[cvRed, rounded corners=2pt] (1.0,0.0) rectangle (1.2,0.45);
    \fill[cvBlue!88, rounded corners=8pt] (-0.8,-0.58) rectangle (0.8,0.78);
    \fill[cvCyan] (-0.4,0.26) circle (0.17);
    \fill[cvCyan] (0.4,0.26) circle (0.17);
    \fill[white] (-0.34,0.33) circle (0.06);
    \fill[white] (0.46,0.33) circle (0.06);
    \draw[cvYellow, line width=2.2pt] (-0.32,-0.22) .. controls (0,-0.46) .. (0.32,-0.22);
  },
  cyber_robot/.pic={
    % shoulders
    \fill[white!92!cvBgMid, rounded corners=6pt] (-0.7,-0.6) -- (-0.8,-0.2) -- (0.8,-0.2) -- (0.7,-0.6) -- cycle;
    \draw[cvNeonCyan, line width=1.5pt] (-0.7,-0.6) -- (-0.8,-0.2) -- (0.8,-0.2) -- (0.7,-0.6) -- cycle;
    % neck
    \fill[cvBgLight] (-0.2,-0.2) rectangle (0.2,0.0);
    \draw[cvNeonCyan, line width=1pt] (-0.2,-0.2) -- (-0.2,0.0) (0.2,-0.2) -- (0.2,0.0);
    % head
    \fill[white, rounded corners=12pt] (-0.6,0.0) rectangle (0.6,0.9);
    \draw[cvNeonBlue, line width=1.8pt, rounded corners=12pt] (-0.6,0.0) rectangle (0.6,0.9);
    % visor/eyes (glowing)
    \fill[cvBgDark, rounded corners=6pt] (-0.45,0.2) rectangle (0.45,0.6);
    \draw[cvNeonCyan, line width=1.2pt, rounded corners=6pt] (-0.45,0.2) rectangle (0.45,0.6);
    \fill[cvNeonCyan] (-0.22,0.4) circle (0.06);
    \fill[cvNeonCyan] (0.22,0.4) circle (0.06);
    \draw[cvNeonCyan, opacity=0.4, line width=1.5pt] (-0.35,0.4) -- (0.35,0.4);
    % antenna
    \draw[cvNeonBlue, line width=2pt] (0,0.9) -- (0,1.25);
    \fill[cvNeonPurple] (0,1.3) circle (0.1);
    \fill[cvNeonPurple, opacity=0.3] (0,1.3) circle (0.22);
  },
  chip/.pic={
    \foreach \x in {-0.32,0,0.32}{
      \fill[cvGrey] (\x,0.55) rectangle ++(0.1,0.14);
      \fill[cvGrey] (\x,-0.69) rectangle ++(0.1,0.14);
    }
    \foreach \y in {-0.32,0,0.32}{
      \fill[cvGrey] (0.55,\y) rectangle ++(0.14,0.1);
      \fill[cvGrey] (-0.69,\y) rectangle ++(0.14,0.1);
    }
    \fill[cvNeonOrange, rounded corners=5pt] (-0.58,-0.58) rectangle (0.58,0.58);
    \node[white, font=\cvHeadLat] at (0,0) {AI};
  },
}

% ---- सहायक मैक्रो: तारा एवं चमक (star & sparkle) ----
\newcommand{\cvstar}[3]{%
  \node[star, star points=5, star point ratio=0.46, fill=#3, scale=#2, inner sep=0pt] at (#1) {};
}
\newcommand{\cvsparkle}[3]{%
  \node[star, star points=4, star point ratio=0.28, fill=#3, scale=#2, inner sep=0pt] at (#1) {};
}

\fi

\begin{titlepage}
\thispagestyle{empty}
\begin{tikzpicture}[remember picture, overlay]
  % ---------- पृष्ठभूमि (dark space gradient) ----------
  \shade[top color=cvBgDark, bottom color=cvBgMid]
        (current page.south west) rectangle (current page.north east);
  
  % Subtle grid pattern
  \draw[step=1.2cm, cvGridColor, line width=0.4pt, opacity=0.45] 
        (current page.south west) grid (current page.north east);

  % Subtle glowing data pathways in the background
  \draw[cvNeonCyan, opacity=0.15, line width=1.5pt] (-10.5, 6.0) .. controls (-3.0, 4.0) and (3.0, -8.0) .. (10.5, -6.0);
  \draw[cvNeonPurple, opacity=0.12, line width=1.2pt] (-10.5, -4.0) .. controls (-4.0, 2.0) and (4.0, 8.0) .. (10.5, 4.0);
  \draw[cvNeonPink, opacity=0.10, line width=1.0pt] (-10.5, 0.0) .. controls (-2.0, -6.0) and (2.0, 6.0) .. (10.5, -2.0);

  \coordinate (C) at (current page.center);

  % Central AI Core Coordinate
  \coordinate (Hub) at ($(C)+(0,-2.0)$);
  
  % Orbits and signals around Hub (geometric orbital shells)
  \draw[cvNeonCyan, opacity=0.06, line width=2pt] (Hub) circle (1.8);
  \draw[cvNeonPurple, opacity=0.10, line width=1.5pt] (Hub) circle (3.0);
  \draw[cvNeonCyan, opacity=0.15, line width=1.2pt] (Hub) circle (5.52);
  
  % Branching Coordinates (symmetrical on the 5.52cm orbit)
  \coordinate (TL) at ($(Hub)+(-4.5, 3.2)$);
  \coordinate (TR) at ($(Hub)+(4.5, 3.2)$);
  \coordinate (BL) at ($(Hub)+(-4.5, -3.2)$);
  \coordinate (BR) at ($(Hub)+(4.5, -3.2)$);
  
  % Glowing neural lines from hub to science branches
  \draw[cvNeonCyan, opacity=0.45, line width=2pt] (Hub) -- (TL);
  \draw[cvNeonGreen, opacity=0.45, line width=2pt] (Hub) -- (TR);
  \draw[cvNeonPink, opacity=0.45, line width=2pt] (Hub) -- (BL);
  \draw[cvNeonYellow, opacity=0.45, line width=2pt] (Hub) -- (BR);
  
  % Sidelinks connecting the science branches
  \draw[white, opacity=0.15, line width=1.2pt, dashed] (TL) -- (BL);
  \draw[white, opacity=0.15, line width=1.2pt, dashed] (TR) -- (BR);
  
  % Inner signals/neurons along paths
  \fill[cvNeonCyan] ($(Hub)!0.45!(TL)$) circle (3.5pt);
  \fill[cvNeonGreen] ($(Hub)!0.55!(TR)$) circle (3.5pt);
  \fill[cvNeonPink] ($(Hub)!0.45!(BL)$) circle (3.5pt);
  \fill[cvNeonYellow] ($(Hub)!0.55!(BR)$) circle (3.5pt);
  
  % Orbital Science Nodes (styled outer rings)
  \fill[cvBgDark, draw=cvNeonPurple, line width=2pt] (TL) circle (1.25);
  \fill[cvBgDark, draw=cvNeonGreen, line width=2pt] (TR) circle (1.25);
  \fill[cvBgDark, draw=cvNeonPink, line width=2pt] (BL) circle (1.25);
  \fill[cvBgDark, draw=cvNeonYellow, line width=2pt] (BR) circle (1.25);
  
  % Icons inside nodes
  \pic at (TL) {atom};
  \pic at (TR) {flask};
  \pic[scale=0.75] at (BL) {dna};
  \pic at (BR) {chip};
  

  
  % Central AI Core Hub (Mascot with glowing halo)
  \fill[cvNeonCyan, opacity=0.08] (Hub) circle (1.4);
  \fill[cvBgDark, draw=cvNeonCyan, line width=2.5pt] (Hub) circle (1.4);
  \pic[scale=1.1] at (Hub) {cyber_robot};

  % Scattered stars & sparkles in background
  \cvstar{$(C)+(-7.8,6.0)$}{0.4}{cvNeonYellow}
  \cvstar{$(C)+(7.8,6.0)$}{0.4}{cvNeonYellow}
  \cvsparkle{$(C)+(-8.5,2.5)$}{0.45}{cvNeonPink}
  \cvsparkle{$(C)+(8.5,2.5)$}{0.45}{cvNeonCyan}
  \cvsparkle{$(C)+(-2.5,-7.5)$}{0.35}{cvNeonPurple}
  \cvsparkle{$(C)+(2.5,-7.5)$}{0.35}{cvNeonGreen}
  \cvstar{$(C)+(8.0,-9.2)$}{0.3}{cvNeonYellow}
  \cvstar{$(C)+(-8.0,-9.2)$}{0.3}{cvNeonYellow}

  % ---------- शीर्षक बैनर (title banner — glassmorphism look) ----------
  \node[fill=cvBgLight!60!cvBgDark, rounded corners=20pt,
        draw=cvNeonCyan, double=cvNeonBlue, double distance=1.5pt, line width=1.2pt,
        inner xsep=24pt, inner ysep=20pt, align=center] (title) at ($(C)+(0,8.7)$) {%
    {\cvTH\fontsize{30}{36}\selectfont\color{white} विज्ञान में}\\[4pt]
    {\cvTH\fontsize{33}{39}\selectfont\color{cvNeonYellow} कृत्रिम बुद्धिमत्ता}\\[10pt]
    {\cvHH\fontsize{15}{18}\selectfont\color{cvNeonCyan} सिद्धांत, अनुप्रयोग एवं भविष्य}\\[8pt]
    {\cvHL\fontsize{14}{16}\selectfont\color{white} Artificial Intelligence in Science}\\[2pt]
    {\cvBL\fontsize{10.5}{13}\selectfont\color{white!70} Principles, Applications and Future}%
  };

  % ---------- टैगलाइन + लेखक (tagline + author pill) ----------
  \node[align=center, text=white] at ($(C)+(0,-9.6)$)
    {\cvHH\fontsize{13}{16}\selectfont एक मज़ेदार, गतिविधि एवं परियोजना आधारित विज्ञान पुस्तक};
    
  \node[fill=cvBgLight!80!cvBgDark, rounded corners=15pt, draw=cvNeonBlue, line width=1.5pt,
        inner xsep=22pt, inner ysep=10pt, text=white] at ($(C)+(0,-11.4)$)
    {\cvTL\fontsize{20}{24}\selectfont Krishna Kumar Godara};
    
  \node[text=cvNeonCyan] at ($(C)+(0,-12.7)$)
    {\cvHL\fontsize{11}{13}\selectfont Department of Physics, IIT Jodhpur ~$\cdot$~ 2026};
\end{tikzpicture}
\end{titlepage}
